\section{\Ks{} and \Lam{} in the default (crossing-matched) sample}
\label{sec:v0}

Figures~\ref{fig:precutKs} and \ref{fig:precutLambda} show the stored topological variables,
i.e.\ after the reconstruction pre-cuts and before any offline cut. Figures~\ref{fig:massKs}
and~\ref{fig:massLambda} show the invariant mass before and after the offline QA cuts of
Section~\ref{sec:qacuts}, after removing duplicates (Section~\ref{sec:dup}). Fits: Gaussian plus
second-order polynomial, binned likelihood, in 0.43--0.57 (\Ks) and 1.09--1.16\,GeV/$c^2$ (\Lam);
the signal yield is the Gaussian integral, S/B is in $\pm3\sigma$. Uncertainties are statistical.

\begin{figure}[H]
\centering
\includegraphics[width=\textwidth]{fig_precut_Ks}
\caption{\Ks{} candidates as stored (pre-cuts only, \nCandKs{} candidates, \nEvents{} events):
DIRA, DCA between the daughters, SV $\chi^2$/ndf, decay length, the smaller of the two daughter
DCA$_{xy}$ to the PV, and mother \pt. Red dashed lines: pre-cut values (DIRA$>0.8$,
DCA$_{12}<0.5$\,cm, $\chi^2$/ndf$<20$). Log scale.}
\label{fig:precutKs}
\end{figure}

\begin{figure}[H]
\centering
\includegraphics[width=\textwidth]{fig_precut_Lambda}
\caption{As Figure~\ref{fig:precutKs} for \Lam+\aLam{} candidates (\nCandLambda{} candidates).}
\label{fig:precutLambda}
\end{figure}

\begin{figure}[H]
\centering
\includegraphics[width=0.9\textwidth]{fig_mass_Ks}
\caption{$\pi^+\pi^-$ invariant mass. Left: pre-cuts only. Right: with offline QA cuts
(\nQAKs{} candidates before, \dupQAKs{} duplicates removed) and the fit (red). The \Ks{} peak is
at $\meanKs\pm\emeanKs$\,MeV, $\massShiftKs$\,MeV below the PDG mass (497.6\,MeV), with
$\sigma=\sigmaKs\pm\esigmaKs$\,MeV, $N_{\rm sig}=\yieldKs\pm\eyieldKs$ and S/B=\SBKs.}
\label{fig:massKs}
\end{figure}

\begin{figure}[H]
\centering
\includegraphics[width=0.9\textwidth]{fig_mass_Lambda}
\caption{$p\pi$ invariant mass (\Lam{} and \aLam{} together). Left: pre-cuts only. Right: with
offline QA cuts and fit: $\mu=\meanLambda\pm\emeanLambda$\,MeV, $\sigma=\sigmaLambda\pm\esigmaLambda$\,MeV,
$N_{\rm sig}=\yieldLambda\pm\eyieldLambda$, S/B=\SBLambda.}
\label{fig:massLambda}
\end{figure}

With the shape fixed to the inclusive fit, the separate yields are
\Lam: $\yieldLambdaP\pm\eyieldLambdaP$ and \aLam: $\yieldLambdaA\pm\eyieldLambdaA$.
The QA-cut signal per 1000 events is \yieldPerKevKs{} (\Ks) and \yieldPerKevLambda{} (\Lam+\aLam);
these are not efficiency corrected and depend strongly on the ad-hoc QA cuts.

\caveat{The \Ks{} mass is $\massShiftKs$\,MeV below the PDG value, i.e.\ about
$-0.6\%$ in mass; with the stated statistical error this is significant. Possible causes
(momentum scale / field map, energy-loss correction, daughter momenta taken at the SV) were not
investigated. The \Lam{} mass agrees with the PDG value within $\approx$1\,MeV.}

\subsection{Yield by daughter type}
\begin{table}[H]
\centering\small
\caption{Signal by daughter-track type (I = INTT, M = MVTX-only, T = TPC-only;
Section~\ref{sec:trackclasses}), default sample with QA cuts, duplicates removed. Gaussian mean
and width fixed to the inclusive fit, polynomial background free.}
\label{tab:pairtype}
\begin{minipage}[t]{0.48\textwidth}\centering
\begin{tabular}{lrrr}
\toprule
\multicolumn{4}{c}{\Ks} \\
types & cand. & signal & \% \\
\midrule
\input{tables/tab_pairtype_Ks.tex}
\bottomrule
\end{tabular}
\end{minipage}\hfill
\begin{minipage}[t]{0.48\textwidth}\centering
\begin{tabular}{lrrr}
\toprule
\multicolumn{4}{c}{\Lam+\aLam} \\
types & cand. & signal & \% \\
\midrule
\input{tables/tab_pairtype_Lambda.tex}
\bottomrule
\end{tabular}
\end{minipage}
\end{table}
Most of the signal has two INTT-matched daughters. Pairs with two TPC-only daughters (TT) carry
a small signal; by construction they come with a crossing-0 PV (Section~\ref{sec:crossing}).
Mass distributions per type: Figures~\ref{fig:pairKs}--\ref{fig:pairLambda}.

\begin{figure}[H]
\centering
\includegraphics[width=0.95\textwidth]{fig_pairtype_Ks}
\caption{$\pi^+\pi^-$ mass by daughter type (default sample, QA cuts, duplicates removed); blue:
fit with mean and width fixed to the inclusive fit. Numbers in Table~\ref{tab:pairtype}.}
\label{fig:pairKs}
\end{figure}
\begin{figure}[H]
\centering
\includegraphics[width=0.95\textwidth]{fig_pairtype_Lambda}
\caption{As Figure~\ref{fig:pairKs} for $p\pi$.}
\label{fig:pairLambda}
\end{figure}
